\documentclass[aspectratio=169,10pt,spanish]{beamer}

\input{../../../_preambulo_beamer.tex}
\input{../../../_comandos_beamer.tex}

\selectlanguage{spanish}

\title{MA1001B - Modelación Estadística para la Toma de Decisiones}
\subtitle{Lección 27: Intervalo de confianza Z para una media (sigma conocida)}
\author{}
\institute{Tecnol\'ogico de Monterrey}
\date{}

\begin{document}

\begin{frame}[plain]
  \vspace{1.2cm}
  {\Large\bfseries MA1001B - Modelación Estadística para la Toma de Decisiones\par}
  \vspace{0.7cm}
  {\large Lección 27: Intervalo de confianza Z para una media ($\sigma$ conocida)\par}
  \vspace{0.7cm}
  {\color{TecRojo}\rule{\textwidth}{1.2pt}}
  \vspace{0.8cm}

  Facultad MA1001B\\[0.3cm]
  Periodo\\[0.3cm]
  Tecnol\'ogico de Monterrey
\end{frame}

\begin{frame}{La pregunta del intervalo con sigma conocida}
  \begin{alertblock}{Pregunta guía}
    ¿Cómo se construye un intervalo al 95\% para una media desconocida
    cuando $\sigma$ se trata como conocida?
  \end{alertblock}

  \vspace{0.6em}
  Al final de esta lección, usted puede:
  \begin{itemize}
    \item calcular $\mathrm{EE}=\sigma/\sqrt{n}$ a partir de una $\sigma$ conocida;
    \item formar $\bar{x}\pm 1.96\cdot\sigma/\sqrt{n}$;
    \item interpretar el intervalo como un rango para $\mu$, no para una botella;
    \item explicar por qué una $\sigma$ desconocida hace de este procedimiento z la herramienta equivocada.
  \end{itemize}

  \vfill
  \scriptsize Todos los volúmenes de llenado usados hoy son sintéticos. No se usan datos reales de empresa.
\end{frame}

\begin{frame}{Una línea de llenado sintética con $\sigma$ conocida}
  \small
  \begin{center}
    \begin{tabular}{lr}
      \toprule
      \textbf{Cantidad} & \textbf{Valor} \\
      \midrule
      Tamaño de muestra $n$ & 36 botellas \\
      $\sigma$ conocida & 10 ml \\
      $\mu$ oculta (no usada por el analista) & 502 ml \\
      Estimación puntual $\bar{x}$ (semilla 42) & 502.738498 ml \\
      \bottomrule
    \end{tabular}
  \end{center}

  \vspace{0.5em}
  \begin{exampleblock}{Error estándar}
    \[
      \mathrm{EE}=\frac{\sigma}{\sqrt{n}}=\frac{10}{\sqrt{36}}=1.666667.
    \]
    El analista conoce $\sigma$ a partir de la capacidad de proceso de largo
    plazo y no conoce $\mu$.
  \end{exampleblock}
\end{frame}

\begin{frame}[fragile]{Paso 1: Media muestral y error estándar}
  \begin{columns}[T]
    \begin{column}{0.42\textwidth}
      \small
      \begin{lstlisting}
se = sigma / np.sqrt(n)
xbar = np.mean(sample)
      \end{lstlisting}

      \begin{block}{Salida verificada}
        $n=36$, $\sigma=10$\\
        $\bar{x}=502.738498$ ml\\
        $\mathrm{EE}=1.666667$ ml
      \end{block}
    \end{column}
    \begin{column}{0.58\textwidth}
      \centering
      \includegraphics[width=\textwidth]{../figuras/z_ci_mean_01_sample_standard_error.png}
      \vspace{0.2em}
      \scriptsize Histograma y banda de EE del Paso 1
    \end{column}
  \end{columns}
\end{frame}

\begin{frame}{Multiplicar $z^{*}$ por el error estándar}
  \small
  Valor crítico de libro al 95\%:
  \[
    z^{*}=1.96.
  \]
  Margen de error e intervalo:
  \[
    E=1.96\times 1.666667=3.266667,
  \]
  \[
    \bar{x}\pm E=[499.471831, 506.005165].
  \]
  Este intervalo con semilla 42 cubre la $\mu$ oculta $=502$. La cobertura
  en una muestra no es un enunciado de probabilidad sobre este intervalo
  realizado.
\end{frame}

\begin{frame}[fragile]{Paso 2: El intervalo Z al 95\%}
  \begin{columns}[T]
    \begin{column}{0.40\textwidth}
      \small
      \begin{lstlisting}
margin = 1.96 * se
lower = xbar - margin
upper = xbar + margin
      \end{lstlisting}

      \begin{block}{Salida verificada}
        $z^{*}=1.96$\\
        Margen $=3.266667$\\
        IC $[499.471831, 506.005165]$
      \end{block}
    \end{column}
    \begin{column}{0.60\textwidth}
      \centering
      \includegraphics[width=\textwidth]{../figuras/z_ci_mean_02_z_interval.png}
      \vspace{0.2em}
      \scriptsize Intervalo z con sigma conocida del Paso 2
    \end{column}
  \end{columns}
\end{frame}

\begin{frame}{Paso 3: Sigma desconocida hace de Z la herramienta equivocada}
  \begin{columns}[T]
    \begin{column}{0.42\textwidth}
      \small
      $s$ muestral $=8.390375\neq 10$.

      \vspace{0.5em}
      Intervalo z-con-$s$ inválido:
      \[
        [499.997642, 505.479354].
      \]

      \vspace{0.4em}
      La anchura cae de $6.533333$ a $5.481711$.

      \vspace{0.5em}
      \textbf{Límite:} conservar $z^{*}=1.96$ después de sustituir $\sigma$
      por $s$ no es un intervalo z. La Lección 28 usa $t$ con
      $\mathrm{gl}=n-1$.
    \end{column}
    \begin{column}{0.58\textwidth}
      \centering
      \includegraphics[width=\textwidth]{../figuras/z_ci_mean_03_unknown_sigma_limit.png}
      \vspace{0.2em}
      \scriptsize Intervalo z válido versus z-con-$s$ inválido del Paso 3
    \end{column}
  \end{columns}
\end{frame}

\begin{frame}{Verificación de conceptos}
  \begin{enumerate}
    \item ¿Por qué $\mathrm{EE}=10/\sqrt{36}$ y no $s/\sqrt{n}$ en esta
      lección?
    \item Calcule el margen al 95\% $1.96\times 1.666667$.
    \item ¿Cubrir $\mu=502$ en esta muestra implica que la siguiente
      muestra cubrirá $\mu$?
    \item ¿Por qué $s=8.390375$ no autoriza conservar $z^{*}=1.96$?
  \end{enumerate}

  \vspace{0.6em}
  \begin{block}{Conexión con la Lección 28}
    La sesión puede continuar directamente con el intervalo $t$ de Student,
    los grados de libertad y el efecto de un solo atípico sobre $s$.
  \end{block}

  \vfill
  \scriptsize Fuente: OpenStax, \textit{Introductory Business Statistics 2e},
  Capítulo 8, Sección 8.1. CC BY-NC-SA.
\end{frame}

\appendix



\end{document}
